\documentclass[11pt]{article}
\usepackage{amsmath,amsthm,amscd,amssymb,mathrsfs,tikz}
\usepackage{latexsym,epsf,epsfig}
\usepackage{sectsty}
\usepackage{hyperref}
\usepackage{graphicx}
\usepackage{booktabs, multirow} % for borders and merged ranges
\usepackage{soul}% for underlines
\usepackage{xcolor,colortbl} % for cell colors
\usepackage{changepage,threeparttable} % for wide tables
\usepackage[left=1.2in, right=1.2in, top=1in, bottom=1in]{geometry} %The above set the parameters adjust the margins. There are many ways to do this, and frankly, what is above is not the best. However, it will work the time being.

\newcommand{\ds}{\displaystyle}
\newcommand{\sU}{\mathscr U}
% Theorem
\theoremstyle{plain}
\newtheorem{theorem}{Theorem}
% Margins
\topmargin=-0.45in
\evensidemargin=0in
\oddsidemargin=0in
\textwidth=6.5in
\textheight=9.0in
\headsep=0.25in

\begin{document}
\begin{proof} If $n \in \mathbb{N}$ and $3n^2+5$ is odd. Then $n$ must be even. \\
	For the sake of contradiction, $n \in \mathbb{N}$ is odd. Suppose: $3n^2+5$ is odd. \\
	If $n$ is odd. Then we can write it in the form $n=2k+1, k \in \mathbb{N}$. This follows: \\
	\begin{align*}
		 & 3n^2+5            \\
		 & = 3(2k+1)^2+5     \\
		 & = 3(4k^2+4k+1)+5  \\
		 & = (12k^2+12k+3)+5 \\
		 & = 12k^2+12k+8
	\end{align*}
	Observe we can simplify this to the following term:
	\begin{align*}
		 & = 12k^2+12k+8     \\
		 & = 4(3k^2+3k+2)    \\
		 & = 2(2(3k^2+3k+2))
	\end{align*}
	Let: $(2(3k^2+3k+2)) = m$ \\
	Note: All values in all prior operations are in $\mathbb{N}$ and all operations preserve closure $m \in \mathbb{N}$
	\[
		2(2(3k^2+3k+2)) = 2m = 3n^2+5
	\]
	Thus $3n^2+5$ is even which is a contradiction thus the original statment "if $n \in \mathbb{N}$ and $3n^2+5$ is odd. Then $n$ must be even." must be true.
\end{proof}
\end{document}
